\documentclass[11pt,a4paper]{article}

\newcommand{\rapporttitre}{Agent}
\newcommand{\rapportversion}{v0.3}
% Mise en page commune à tous les rapports du projet.
% Chaque rapport définit \rapporttitre et \rapportversion avant de charger ce fichier.

\usepackage{fontspec}
\usepackage[french]{babel}
\usepackage[a4paper,margin=2.2cm,bottom=2.6cm]{geometry}
\usepackage{graphicx}
\usepackage{float}
\usepackage{booktabs}
\usepackage{tabularx}
\usepackage{array}
\usepackage[table]{xcolor}
\usepackage{caption}
\usepackage{enumitem}
\usepackage{titlesec}
\usepackage{fancyhdr}
\usepackage{hyperref}
\usepackage{xurl}

\setmainfont{texgyreheros}[
  Extension=.otf,
  UprightFont=*-regular,
  BoldFont=*-bold,
  ItalicFont=*-italic,
  BoldItalicFont=*-bolditalic]
\definecolor{accent}{HTML}{2A78D6}
\definecolor{ink2}{HTML}{52514E}
\definecolor{fond}{HTML}{EEF4FC}
\hypersetup{colorlinks=true, linkcolor=accent, urlcolor=accent}
\graphicspath{{figures/}}

\titleformat{\section}{\Large\bfseries\color{accent}}{\thesection.}{0.6em}{}[\vspace{-0.4em}\color{accent}\rule{\linewidth}{0.8pt}]
\titleformat{\subsection}{\large\bfseries}{\thesubsection}{0.6em}{}
\titlespacing*{\section}{0pt}{1.6em}{0.9em}

\captionsetup{font={small,color=ink2}, labelfont=bf, justification=raggedright, singlelinecheck=false}
\setlist{itemsep=0.25em, topsep=0.4em}
\setlength{\parskip}{0.45em}
\renewcommand{\arraystretch}{1.25}
\newcolumntype{Y}{>{\raggedright\arraybackslash}X}

\pagestyle{fancy}
\fancyhf{}
\fancyhead[L]{\small\color{ink2} complaint-triage -- \rapporttitre}
\fancyhead[R]{\small\color{ink2} \rapportversion}
\fancyfoot[R]{\small\color{ink2}\thepage}
\renewcommand{\headrulewidth}{0.3pt}

% \pagedegarde{sous-titre}{date}{ligne de version}
\newcommand{\pagedegarde}[3]{%
  \hypersetup{pageanchor=false}
  \begin{titlepage}
    \setlength{\parindent}{0pt}
    \vspace*{5cm}
    {\color{accent}\bfseries\large COMPLAINT-TRIAGE}\par\vspace{0.6cm}
    {\fontsize{30}{36}\selectfont\bfseries \rapporttitre}\par\vspace{0.5cm}
    {\Large\color{ink2} #1}\par
    \vfill
    {\color{accent}\rule{\linewidth}{1.2pt}}\par\vspace{0.3cm}
    {\color{ink2} Nathan Mang\\ #2\\ #3}
    \vspace{2cm}
  \end{titlepage}
  \hypersetup{pageanchor=true}}

% Encadré pour les points à retenir.
\newcommand{\aretenir}[1]{%
  \par\noindent\colorbox{fond}{\parbox{\dimexpr\linewidth-2\fboxsep}{\small #1}}\par}

\usepackage{microtype}

\begin{document}

\pagedegarde{Un agent qui annonce son plan, utilise des outils encadrés et refuse en cas de doute}{Version de travail, octobre 2026}{Étapes 1 et 2 sur 5 rédigées}

\setcounter{tocdepth}{1}
\tableofcontents
\newpage

\section{Objectif et plan}

\subsection{Ce que doit faire l'agent}

Le cadrage décrit un agent qui lit une réclamation, annonce son plan, puis agit avec des outils encadrés : calculer la priorité d'un dossier, chercher la règle applicable, interroger la base, mettre à jour un statut, alerter un responsable. S'il manque d'éléments, ou si la demande est ambiguë, il refuse et l'explique (besoin B4). L'outil change l'ordre de traitement et prépare le travail ; il ne décide jamais de l'issue d'un dossier.

Les deux versions précédentes fournissent ses outils de raisonnement : la règle de priorisation retenue en v0.1, et le RAG de la v0.2, qui répond en citant les textes ou refuse. La v0.3 ajoute la base de données, les outils d'action et l'agent qui les orchestre. L'évaluation formelle de l'agent, sur 20 à 25 scénarios écrits à l'avance, est prévue en v0.4.

\subsection{Étapes}

\begin{table}[H]
\centering\small
\begin{tabularx}{\linewidth}{@{}l Y@{}}
\toprule
Étape & Contenu \\
\midrule
1 & Base de données : comptes au moindre privilège, schéma, file de réclamations \\
2 & Outils de l'agent, sans modèle de langage : entrées validées, requêtes paramétrées \\
3 & Le graphe de l'agent : plan, validation du plan par le code, exécution, réponse ou refus \\
4 & Premiers scénarios de démonstration \\
5 & Bilan, rapport et note de décision \\
\bottomrule
\end{tabularx}
\caption{Étapes de la v0.3}
\end{table}

\section{La base de données}

\subsection{Pourquoi PostgreSQL, installé sur la machine}

Le cadrage a retenu PostgreSQL, courant en entreprise, plutôt que SQLite, peu représentatif d'un environnement de production. Il est installé directement sur le poste (version 18.6), en service Windows à démarrage manuel : il ne consomme rien quand le projet n'est pas utilisé. Docker, plus lourd sur une machine de 8~Go, arrivera en v0.5 avec l'industrialisation.

\subsection{Trois comptes, au moindre privilège}

Un agent qui agit sur une base doit être limité par la base elle-même, pas seulement par son code. Si le code ou le modèle de langage se trompait, les droits empêcheraient l'erreur.

\begin{table}[H]
\centering\small
\begin{tabularx}{\linewidth}{@{}l Y@{}}
\toprule
Compte & Droits \\
\midrule
\texttt{ct\_owner} & Crée les tables et charge les données. Jamais utilisé par l'agent. \\
\texttt{ct\_agent\_read} & Lecture seule des réclamations et du journal ; transactions en lecture seule par défaut. \\
\texttt{ct\_agent\_write} & Modifie uniquement les deux colonnes de statut d'un dossier, et ajoute des lignes au journal. \\
\bottomrule
\end{tabularx}
\caption{Les comptes de la base}
\end{table}

Les deux comptes de l'agent ont aussi une limite de 5 secondes par requête, pour qu'une requête ne puisse jamais bloquer la base. Aucun ne peut supprimer quoi que ce soit.

Les droits sont vérifiés par des tests qui tentent réellement les opérations interdites sur la base : écrire avec le compte de lecture, y compris après avoir désactivé le mode lecture seule ; modifier la note de priorité ou la banque d'un dossier ; supprimer une réclamation ou une ligne du journal ; écrire un statut inconnu. PostgreSQL refuse chacune d'elles. Ces tests demandent la base : ils tournent sur le poste et sont sautés en CI.

\subsection{La mise en place et les secrets}

Un script crée les comptes, la base, le schéma et les droits. Il demande le mot de passe du super-utilisateur en saisie masquée et ne l'enregistre nulle part. Les mots de passe des trois comptes sont tirés au hasard (32 caractères) et écrits dans le fichier \texttt{.env}, sans jamais être affichés ; le script vérifie d'abord que Git ignore bien ce fichier. Relancer le script change les mots de passe, ce qui sert en cas de fuite. Les ordres SQL sont composés avec les outils de psycopg, qui échappent les identifiants et les valeurs ; un test vérifie qu'un mot de passe contenant « \texttt{'; DROP ROLE} » reste un simple texte.

La première exécution a échoué sur un défaut propre à Windows. psycopg transmet le mot de passe à libpq, la bibliothèque de connexion de PostgreSQL, sans guillemets s'il ne contient pas d'espace. Or, sous Windows, libpq prend l'octet \texttt{0xA0} pour un espace, et cet octet fait partie du codage UTF-8 de « à ». Un mot de passe contenant ce caractère, fréquent sur un clavier français, était donc coupé en deux. Je l'ai établi en reproduisant l'erreur avec des mots de passe factices, sans jamais connaître le vrai. Le mot de passe administrateur passe désormais par la variable \texttt{PGPASSWORD}, lue telle quelle, propre au processus et effacée à la fin.

\subsection{Le schéma et la file}

La table des réclamations ne contient que ce qu'une banque connaît à l'arrivée d'une réclamation : produit, motif, établissement, État, canal, étiquettes et récit s'il existe. Elle ne contient ni l'issue, qui est la cible du modèle, ni le code postal, dont l'agent n'a pas besoin : c'est le principe de minimisation des données. S'y ajoutent la note de priorité et un statut, limité à quatre valeurs (nouveau, en cours, traité, escaladé). Un journal enregistre chaque action de l'agent, avec son auteur et sa date.

La file contient les 140\,512 réclamations reçues du 1\textsuperscript{er} janvier au 16 octobre 2026, auprès de 691 établissements. Leur note est celle de la règle retenue en v0.1, apprise sur les issues connues jusqu'au 31 décembre 2025, comme elle le serait en production. Le chargement passe par une table temporaire et la commande \texttt{COPY}, puis insère ou met à jour les réclamations d'un bloc ; le relancer ne modifie pas les statuts déjà changés.

\section{Les outils de l'agent}

\subsection{Le principe : l'agent choisit, le code exécute}

L'agent ne peut jamais écrire de SQL ni lancer une commande. Il choisit un outil dans une liste fermée et lui donne des valeurs. Chaque outil décrit ses entrées par un modèle pydantic : types, bornes et valeurs permises. Les entrées sont contrôlées avant toute action ; un champ inconnu, une valeur hors bornes ou un type inattendu (« 123 » en texte au lieu d'un nombre) sont refusés, avec un message qui dit quoi corriger, plutôt que convertis en silence. Les requêtes sont écrites dans le code, et les valeurs y passent toujours en paramètres.

\begin{table}[H]
\centering\small
\begin{tabularx}{\linewidth}{@{}l Y l@{}}
\toprule
Outil & Rôle & Compte \\
\midrule
\texttt{find\_company} & Retrouver le nom exact d'un établissement à partir d'une partie du nom & lecture \\
\texttt{get\_complaint} & Lire une réclamation & lecture \\
\texttt{priority\_context} & Rang du dossier dans sa file de la semaine (même banque, même produit), et s'il fait partie des 10~\% traités en premier & lecture \\
\texttt{list\_priority\_queue} & Dossiers les plus prioritaires d'un établissement & lecture \\
\texttt{search\_regulations} & Question de réglementation, par le RAG de la v0.2 & aucun \\
\texttt{update\_status} & Changer le statut d'un dossier, avec une raison & écriture \\
\texttt{raise\_alert} & Alerter un responsable sur un dossier & écriture \\
\bottomrule
\end{tabularx}
\caption{Les sept outils}
\end{table}

\subsection{Les garde-fous de chaque outil}

\begin{itemize}
  \item \textbf{Les noms d'établissement doivent être exacts.} Un nom approché (« Banque fictive » au lieu du nom complet) est refusé, avec l'invitation à utiliser \texttt{find\_company}. L'agent ne peut pas agir sur une banque qu'il a devinée. Dans la recherche par nom, les caractères spéciaux de \texttt{LIKE} sont neutralisés : « \% » cherche un vrai « \% », et non n'importe quel texte.
  \item \textbf{Un dossier ne recule jamais.} \texttt{update\_status} n'accepte que « en cours », « traité » ou « escaladé » : l'agent ne peut pas remettre un dossier à « nouveau ». Une raison est obligatoire. Redonner le même statut est refusé, ce qui évite des écritures en double si l'agent répète une action.
  \item \textbf{Chaque écriture est tracée dans la même transaction.} Le changement de statut et sa ligne au journal (statut avant, après, raison) sont validés ensemble ou annulés ensemble. Pour une alerte, l'inscription au journal n'a lieu que si l'envoi réussit ; un envoi qui échoue ne laisse aucune trace d'une alerte qui n'est pas partie.
  \item \textbf{Les alertes ne sortent pas du poste.} Elles sont envoyées en JSON à un webhook. L'adresse doit être en https, ou en http vers la machine locale seulement. Pendant le développement, un petit récepteur, qui n'écoute que sur 127.0.0.1, les enregistre dans un fichier ; un salon Discord ou Slack pourra le remplacer sans changer l'agent.
  \item \textbf{Des récits tronqués.} Un récit est renvoyé limité à 2\,000 caractères, pour ne pas surcharger le modèle de langage.
\end{itemize}

\subsection{Les tests}

Deux niveaux de tests. Sans base, 20 tests vérifient que les entrées invalides sont refusées avant toute ouverture de la base : identifiant négatif ou donné en texte, champ inconnu nommé « sql », limite de 500 dossiers, produit inexistant, retour au statut « nouveau », alerte de gravité inventée, adresse de webhook non sûre. L'objet base utilisé échoue s'il est ouvert, ce qui prouve que le refus vient de la validation. Avec la base, 10 tests utilisent cinq réclamations d'un établissement fictif, avec des identifiants hors de portée des vrais : ordre de la file, rang dans la semaine, refus d'un nom approché, journal d'un changement de statut, refus d'un statut redonné, arrivée d'une alerte au récepteur, et absence de trace quand l'envoi échoue.

\end{document}
